% !TEX program = xelatex
% Hefu Shan — academic/research and engineering CV (November 2026).
% Based on the MIT-licensed billryan/resume class; no font files included.
% Compile with XeLaTeX; keep this .tex and resume.cls in the same folder.
\documentclass{resume}
\hypersetup{pdftitle={Hefu Shan - Curriculum Vitae},pdfauthor={Hefu Shan},pdfsubject={Research, Machine Learning and Software Engineering}}
\begin{document}
\pagenumbering{gobble}
\name{Hefu Shan}
\basicInfo{\faMapMarker\ Tokyo, Japan / Beijing, China \quad\textperiodcentered\quad
  \faEnvelopeO\ \href{mailto:shanhefu@gmail.com}{shanhefu@gmail.com} \\[1pt]
  \faPhone\ (+81) 080-7417-1012 \quad\textperiodcentered\quad \faPhone\ +86 18610916979}
\secondInfo{\faGithub\ \href{https://github.com/shengruduzhou}{github.com/shengruduzhou} \quad\textperiodcentered\quad
  \faGlobe\ \href{https://shengruduzhou.github.io/}{shengruduzhou.github.io} \quad\textperiodcentered\quad
  \href{https://shanhefu-portfolio.vercel.app/}{Portfolio} \quad\textperiodcentered\quad
  \faLinkedin\ \href{https://linkedin.com/in/hefu-shan-054b24361/}{LinkedIn}}
\vspace{-1.9pt}

\sectionwithicon{\faGraduationCap}{Education}
\datedsubsection{Sophia University, Tokyo, Japan}{Apr. 2025 -- Present}
\role{Master of Science in Artificial Intelligence (Computer Vision), expected Mar. 2027}
\vspace{4pt}
\datedsubsection{Teikyo University, Utsunomiya, Japan}{Apr. 2021 -- Mar. 2025}
\role{Bachelor of Science in Information and Electronics Engineering}

\sectionwithicon{\faLightbulbO}{Research Interests}
Self-Supervised Learning; Representation Learning; Panoramic Images; Spherical Geometry.

\sectionwithicon{\faFlask}{Research Experience}
\datedsubsection{PanoMAE --- Self-Supervised Learning for Panoramic Vision}{Jun. 2025 -- Present}
\role{Graduate Research, Sophia University \quad\textperiodcentered\quad \href{https://github.com/islab-sophia-edu/2026-SHAN-HEFU/issues}{Research Notes}}
\begin{itemize}
  \item Developed a masked-autoencoding framework for panoramic images using tangent-plane sampling, spherical positional encoding, and geometry-aware attention.
  \item Evaluated learned representations on classification, semantic segmentation, depth and surface-normal estimation; studied geometric ablations and rotation robustness.
\end{itemize}

\sectionwithicon{\faCode}{Selected Projects}
\cvproj{\href{https://github.com/shengruduzhou/QuantAgent}{QuantAgent}}{2026 -- Present}
 {Python, Qlib, Polars, FastAPI, React}{Built a point-in-time A-share research system with walk-forward evaluation, cost-aware backtesting and portfolio risk controls.}
\cvproj{\href{https://github.com/shengruduzhou/FormulaForge}{FormulaForge}}{2026 -- Present}
 {TypeScript, React, Node.js, LaTeX, LLM APIs}{Developed a formula-analysis tool combining LaTeX parsing, screenshot/OCR input and structured multimodal explanations.}
\cvproj{\href{https://github.com/shengruduzhou/Feedback_System}{Online Course Feedback System}}{Apr. 2024 -- May 2025}
 {Bachelor's Thesis \quad\textperiodcentered\quad Flask, React, MySQL, ARIMA}{Developed a web application for analyzing and visualizing student concentration and classroom feedback.}
\cvproj{T-net Classroom Question System}{2023}
 {Group Project \quad\textperiodcentered\quad Java, Spring Boot, JavaScript, CSS, MySQL}{Built a web-based quiz system supporting true/false questions and automatic score calculation.}
\cvproj{Blog and Portfolio}{May 2025 -- Present}
 {Jekyll, React, TypeScript, CSS, GSAP}{Created a \href{https://shengruduzhou.github.io/}{technical blog} and an \href{https://shanhefu-portfolio.vercel.app/}{interactive portfolio} to document projects and research.}

\sectionwithicon{\faCogs}{Technical Skills}
\textbf{Programming:} Python, Java, JavaScript/TypeScript, SQL\\
\textbf{Machine Learning:} PyTorch, OpenCV, NumPy, pandas, scikit-learn, Vision Transformers\\
\textbf{Software and Data:} React, FastAPI, Flask, Spring Boot, Node.js, Polars, MySQL, PostgreSQL\\
\textbf{Tools:} Linux, Git, Docker, CUDA, Conda, \LaTeX

\sectionwithicon{\faTrophy}{Honors and Additional Information}
\textbf{Scholarship:} Japan Student Services Organization (JASSO), recipient (2021).\\
\textbf{Recognition:} Rotary Yoneyama Memorial Scholarship, final interview nominee (2023).\\
\textbf{Languages:} Chinese (native); Japanese (fluent); English (intermediate, TOEIC 610).
\end{document}